% !TeX root = ../../../main.tex
\section{Introduction}

Although short-term survival rates for transplanted kidneys have improved in recent years~\cite{saritas_Kidney_2021}, detection and monitoring of long-term graft health remains a significant challenge~\cite{poggio_Longterm_2021}. The main cause of graft loss is antibody-mediated rejection (AMR), which is characterized by the presence of inflammation and fibrosis in the renal parenchyma~\cite{loupy_Outcome_2009,nankivell_Rejection_2010}. The detection of AMR is difficult using noninvasive techniques, and the current gold standard for diagnosing AMR is a needle biopsy~\cite{alexandre_AntibodyMediated_2018}. However, biopsies are invasive, painful, expensive, and can cause complications such as bleeding, infection, and graft loss~\cite{schinstock_Value_2017}. Thus, there is a significant clinical need for noninvasive techniques that can detect and monitor inflammation and fibrosis in the renal parenchyma.

One of the hallmarks of AMR is the presence of parenchymal inflammation, an acute reaction by the immune system to injuries~\cite{loupy_Banff_2020,mannon_Inflammation_2010,alexandre_AntibodyMediated_2018}. The risk of renal graft loss due to AMR is exacerbated when chronic immune responses such as fibrosis are also present in the kidney~\cite{venner_Relationships_2016, mannon_Inflammation_2010, saritas_Kidney_2021}. It should be noted, though, that inflammation and fibrosis can also occur from other causes such as infections or toxicity effects from immunosuppresants including calcineurin inhibitors~\cite{naesens_Proteinuria_2016, callemeyn_Revisiting_2021}. However, the simultaneous presence of both pathologies is a strong indicator of AMR. Thus, it is essential to differentiate between the presence of inflammation alone, fibrosis alone, and the combination of both pathologies in the renal parenchyma.

Current clinical standards require regular surveillance of renal allograft health by regularly collecting serum and urine samples~\cite{naesens_Proteinuria_2016}, taking multiple measurements of nonspecific biomarkers such as blood urea nitrogen (BUN) and serum creatinine concentrations, and performing needle biopsies when their use is clinically indicated~\cite{loupy_Banff_2020, schinstock_Value_2017}. However, blood and urine tests are neither specific for inflammation nor fibrosis in the renal parenchyma. Thus, there is a significant clinical need for noninvasive techniques that can specifically detect and differentiate inflammation and fibrosis in kidneys.

One potential technique to meet this need is ultrasound elastography~\cite{urban_Novel_2021}. Ultrasound elastography is a noninvasive imaging technique that measures the mechanical properties of tissues by mechanically exciting them and observing their motion over time. One common approach is acoustic radiation force impulse (ARFI) imaging~\cite{nightingale_feasibility_2001,nightingale_Acoustic_2002}, where acoustic radiation force (ARF) excitations are emitted in the form of ultrasonic waveforms that are up to tens of milliseconds long, and the resulting displacement at the region of excitation (RoE) is tracked over time using a train of shorter ultrasonic pulses. Another approach is shear wave elasticity imaging (SWEI)~\cite{sarvazyan_Shear_1998,nightingale_Shearwave_2003}, where a shear wave is generated by an ARF excitation, and the propagating shear wave is tracked over time across multiple locations away from the RoE.

Although ARFI and SWEI have been used to characterize multiple pathologies~\cite{tsai_Sonographic_2013, stock_ARFIbased_2011, magalhaes_Diagnostic_2017, frulio_Evaluation_2013, barr_Update_2020}, both techniques present challenges in differentiating inflammation and fibrosis in the renal parenchyma. Because the magnitude and distribution of force amplitude applied onto tissue from an ARF excitation is unknown, ARFI is ultimately a qualitative technique. As such, ARFI peak displacements cannot be directly compared between images or patients. Although SWEI is quantitative in that it can quantify shear wave speed (and shear elastic modulus, given additional assumptions) across images, SWEI measurements assume that tissues are isotropic and homogeneous between the RoE and displacement-tracking site~\cite{sarvazyan_Shear_1998}. This assumption breaks down in the renal parenchyma, where tissues are heterogeneous and anisotropic at sub-millimeter scales~\cite{urban_Novel_2021, liu_Effect_2017}. Furthermore, since displacement measurements for SWEI take several milliseconds, measurements can be corrupted by the reflection and dissipation of shear waves. Thus, traditional approaches to ultrasound elastography face significant challenges when their use is attempted in the renal parenchyma~\cite{orlacchio_Kidney_2014, urban_Novel_2021}.

One strategy to overcome this challenge is to develop techniques that combine aspects of ARFI and SWEI. One example is Viscoelastic Response (VisR) elastography~\cite{selzo_Quantitative_2016}, a technique that has been demonstrated to characterize the health of transplanted kidneys in humans \emph{in vivo}~\cite{hossain_Evaluating_2018, yokoyama_vivo_2022}. VisR applies two ARF excitations in short succession, tracks induced displacements at the RoE over time, and fits a one-dimensional mass-spring-damper model to the data. The model parameters allow for relative elasticity (RE) and relative viscosity (RV) measurements, which relate to shear elastic and viscous moduli~\cite{hossain_Viscoelastic_2020}, to be assessed in isotropic materials as well as ~\emph{in silico} anisotropic materials~\cite{muhtadi_VisR_2024}. Although RE and RV also relate to a force amplitude term of unknown magnitude, ratios of their measurements can be calculated between two different transducer orientations or anatomical features to derive semi-quantitative degrees of anisotropy (DoA) or regional ratio (RR) metrics, respectively. Using DoA and RR, VisR was shown to differentiate renal allografts that only clinically indicated fibrosis from those that presented both fibrosis and inflammation~\cite{yokoyama_vivo_2022}. However, those previous works did not investigate VisR's abilities to differentiate inflammation alone or to distinguish kidneys that have inflammation, fibrosis, and both pathologies from uninjured kidneys. It is critical to assess whether VisR can identify those different pathologies to establish its clinical utility for identifying AMR, thereby making it possible for clinicians to intervene earlier and prevent graft loss. Furthermore, a question remains of whether the semiquantitative metrics derived from VisR suffice for identifying AMR-relevant lesions, or whether a truly quantitative elastography technique is required for such an application.

A potential pathway to improve the identification of inflammation, fibrosis, and their combination is double-profile intersection (DoPIo) elastography~\cite{yokoyama_DoubleProfile_2025, yokoyama_Doubleprofile_2019}. DoPIo is a technique that applies a single ARF excitation, tracks induced displacements at the RoE, and then beamforms the same channel data into two displacement profiles that encode the same underlying tissue mechanical response~\cite{yokoyama_Quantitative_2022}. The paired displacement measurements are colocated but created using different tracking apertures (e.g. using identical channel data but beamforming using two different F-numbers) such that one tracking beam is wider than the other. As a result, the two displacement profiles encode the same scatterers' motion while applying different spatial weights to the echoes of displaced scatterers. The two profiles are overlapped, and the time when they intersect (i.e. the time-intersect) is measured as a common measure of scatterers' shear rate. Finally, times-intersect are converted into shear elasticity estimates using a linear model empirically derived using finite element model simulations of homogeneous, linearly elastic materials undergoing ARF excitations. Using this approach, DoPIo has accurately quantified shear elasticity \emph{in vitro} in calibrated phantoms and demonstrated comparable elasticity measurements in \emph{ex vivo} tissue to SWEI~\cite{yokoyama_Doubleprofile_2019,yokoyama_Assessing_2021}. DoPIo has also been shown to exhibit sensitivity to elastic anisotropy in transversely isotropic (TI) materials, including a dependence on longitudinal Young's moduli \emph{in silico}~\cite{muhtadi_Double_2025} as well as to differentiate elastic features in canine kidneys \emph{ex vivo}~\cite{yokoyama_Assessing_2021}.

With the recent developments of DoPIo, we aimed to investigate whether DoPIo could distinguish pathological changes of elasticity in renal parenchyma \emph{in vivo}. We hypothesized that DoPIo could differentiate between inflamed and fibrotic kidneys, and those with both pathologies. We also aimed to compare the performance of DoPIo to VisR in differentiating these pathologies, and to determine whether DoPIo could act as a better diagnostic marker for kidneys with acute, chronic, and/or combined injuries. To perform this comparison, this work describes the implementation of VisR and DoPIo using a preclinical model of ischemic reperfusion injuries (IRI) to induced inflammation, fibrosis, and the combination of inflammation and fibrosis as seen in AMR.